\documentclass[12pt,a4paper]{article}

\usepackage[a4paper,margin=2cm,headheight=15pt,footskip=22pt]{geometry}
\usepackage[T2A]{fontenc}
\usepackage[utf8]{inputenc}
\usepackage[russian]{babel}
\usepackage{amsmath,amssymb,mathtools}
\usepackage{array,booktabs,longtable}
\usepackage{xcolor}
\usepackage{hyperref}
\usepackage{tikz}
\usepackage{fancyhdr}

\definecolor{teal}{RGB}{0,128,128}
\hypersetup{
  colorlinks=true,
  linkcolor=blue!55!black
}

\setlength{\parindent}{0pt}
\setlength{\parskip}{0.65em}
\renewcommand{\arraystretch}{1.25}

\pagestyle{fancy}
\fancyhf{}
\fancyfoot[C]{\small\color{gray!80!black}Лёвин Артём Александрович эксклюзивно для Ксении Васильченко}
\fancyfoot[R]{\small\thepage}
\renewcommand{\headrulewidth}{0pt}
\renewcommand{\footrulewidth}{0pt}

\newcolumntype{P}[1]{>{\raggedright\arraybackslash}p{#1}}
\newcommand{\ErrorLabel}{\textcolor{red}{Ошибка:}}
\newcommand{\KeyIdea}{\textcolor{blue!60!black}{Ключевая идея:}}
\newcommand{\Outcome}{\textcolor{teal!60!black}{Итог:}}

\begin{document}

\begin{titlepage}
\thispagestyle{fancy}
\centering
\vspace*{0.24\textheight}
{\Huge\bfseries Ревью домашней работы\par}
\vspace{1.8cm}
{\large Автор: Лёвин Артём Александрович\par}
\vspace{0.7cm}
{\large 20 сентября 2026 года\par}
\vfill
\begin{tikzpicture}[x=1cm,y=1cm]
  \draw[blue!35!black,line width=0.8pt]
  (0,0) .. controls (1.6,0.55) and (3.4,-0.55) .. (5,0)
        .. controls (6.4,0.45) and (7.8,-0.25) .. (9,0.1);
\end{tikzpicture}
\vspace{0.6cm}
\end{titlepage}

\newpage
\section*{Сводная таблица}
\small
\setlength{\tabcolsep}{3pt}
\begin{longtable}{P{0.07\linewidth} P{0.18\linewidth} P{0.28\linewidth} P{0.18\linewidth} P{0.18\linewidth}}
\toprule
\textbf{№ задания} & \textbf{Тема} & \textbf{Суть ошибки} & \textbf{Ваш ответ} & \textbf{Правильный ответ} \\
\midrule
\endfirsthead
\toprule
\textbf{№ задания} & \textbf{Тема} & \textbf{Суть ошибки} & \textbf{Ваш ответ} & \textbf{Правильный ответ} \\
\midrule
\endhead
\hyperref[task:1g]{1г} & Деление смешанных чисел & Неверно переведены оба смешанных числа в неправильные дроби. & $\dfrac{42}{4}$ & $1\dfrac{1}{2}$ \\
\addlinespace
\hyperref[task:2b]{2б} & Признаки делимости на 3 и 4 & Среди подходящих вариантов оставлены числа, которые не делятся на 3; условия не пересечены. & $1224;\ 1204$\newline $1264;\ 1284$ & цифры $2$ и $8$ \\
\addlinespace
\hyperref[task:3]{3} & Разложение на простые множители & Вместо числа $72$ разложено число $70$. & для $72$ использовано $70$; $90$ разложено верно & $72=2^3\cdot3^2$; $90=2\cdot3^2\cdot5$ \\
\addlinespace
\hyperref[task:5]{5} & НОД и сокращение дроби & В НОД включён множитель $7$, которого нет в разложении $414$, и пропущен второй общий множитель $3$; дробь не сокращена. & $\text{НОД}=42$; дробь не сокращена & $\text{НОД}=18$; $\dfrac{23}{28}$ \\
\bottomrule
\end{longtable}
\normalsize

\newpage
\thispagestyle{fancy}
\vspace*{\fill}
\begin{center}
{\Huge\bfseries Разбор ошибок}
\end{center}
\vspace*{\fill}

\newpage
\phantomsection\label{task:1g}
\section*{Задание 1г}

\subsection*{Текст задания}
Вычислите:
\[
2\frac{1}{4}:1\frac{1}{2}.
\]

\subsection*{Ваше решение}
В работе записан переход
\[
\frac{21}{4}:\frac{1}{2}
=\frac{21}{4}\cdot\frac{2}{1}
=\frac{42}{4}.
\]

\subsection*{Ваш ответ}
Последняя записанная величина --- $\dfrac{42}{4}$.

\subsection*{Ошибка}
\ErrorLabel\ неверный перевод смешанных чисел в неправильные дроби.

\subsection*{Где именно возникает ошибка}
До перевода смешанных чисел операция переписана по смыслу верно: нужно разделить первое смешанное число на второе. Первый достоверно неверный переход возникает при замене
\[
2\frac14\longrightarrow\frac{21}{4}.
\]
Одновременно второе смешанное число заменено так:
\[
1\frac12\longrightarrow\frac12,
\]
то есть его целая часть потеряна.

\subsection*{Почему этот шаг неверен}
Смешанное число не переводят в неправильную дробь приписыванием целой части к числителю и не переводят удалением целой части. Для числа вида
\[
a\frac{b}{c}
\]
числитель неправильной дроби получают так: целую часть умножают на знаменатель и прибавляют числитель дробной части. Знаменатель остаётся прежним.

Поэтому для первого числа нужно вычислить: два умножить на четыре и прибавить один. Получается девять, значит
\[
2\frac14=\frac94.
\]
Для второго числа: один умножить на два и прибавить один. Получается три, значит
\[
1\frac12=\frac32.
\]

Запись $\dfrac{21}{4}$ задаёт число, равное пяти целым одной четверти, поэтому она существенно больше исходного числа две целых одна четверть. Запись $\dfrac12$ задаёт одну вторую и, наоборот, существенно меньше исходного числа одна целая одна вторая. После такой замены дальнейшее правило деления дробей применяется уже не к тем числам, которые даны в условии.

\subsection*{Ключевая идея}
\KeyIdea\ сначала каждое смешанное число нужно правильно перевести в неправильную дробь, а затем деление дробей заменить умножением на дробь, обратную делителю.

\subsection*{Исправление от последнего правильного шага}
Возвращаемся к исходным смешанным числам и выполняем правильный перевод:
\[
2\frac14:1\frac12
=\frac94:\frac32
=\frac94\cdot\frac23.
\]
Теперь можно сократить общий множитель: девять и три сокращаются в три раза, а два и четыре --- в два раза. Получаем
\[
\frac94\cdot\frac23=\frac32=1\frac12.
\]

\subsection*{Полное правильное решение}
Переведём первое смешанное число:
\[
2\frac14=\frac{2\cdot4+1}{4}=\frac94.
\]
Переведём второе смешанное число:
\[
1\frac12=\frac{1\cdot2+1}{2}=\frac32.
\]
Деление на дробь заменяем умножением на обратную дробь:
\[
\frac94:\frac32=\frac94\cdot\frac23.
\]
Перемножаем после сокращения:
\[
\frac94\cdot\frac23
=\frac{9\cdot2}{4\cdot3}
=\frac{18}{12}
=\frac32
=1\frac12.
\]

\subsection*{Проверка}
Если частное равно одной целой одной второй, то при умножении его на делитель должно получиться делимое:
\[
1\frac12\cdot1\frac12
=\frac32\cdot\frac32
=\frac94
=2\frac14.
\]
Получено исходное делимое, значит вычисление согласуется с условием.

\subsection*{Итог}
\Outcome\ ошибка возникла не в правиле деления дробей, а раньше --- при переводе обоих смешанных чисел. Правильное значение выражения равно $1\dfrac12$.

\subsection*{Задача на закрепление}
Вычислите
\[
3\frac12:1\frac34.
\]

\newpage
\phantomsection\label{task:2b}
\section*{Задание 2б}

\subsection*{Текст задания}
В числе $12\square4$ вместо квадрата поставьте все цифры, при которых число одновременно делится на $3$ и на $4$.

\subsection*{Ваше решение}
В читаемой части работы перечислены варианты
\[
1224;\qquad 1204;\qquad 1264;\qquad 1284.
\]

\subsection*{Ваш ответ}
Отдельный итоговый список цифр не записан; по записи в качестве вариантов оставлены цифры $2$, $0$, $6$ и $8$.

\subsection*{Ошибка}
\ErrorLabel\ не выполнено одновременное применение двух признаков делимости.

\subsection*{Где именно возникает ошибка}
Первый перечисленный вариант $1224$ действительно подходит: последние две цифры образуют число $24$, которое делится на $4$, а сумма цифр равна девяти, поэтому число делится на $3$.

Первый достоверно неверный вариант --- $1204$. Последние две цифры образуют число $04$, то есть $4$, поэтому условие делимости на $4$ выполнено. Но сумма цифр равна
\[
1+2+0+4=7,
\]
а семь на три не делится. Значит $1204$ нельзя оставлять среди ответов. Аналогично для $1264$ сумма цифр равна тринадцати, поэтому это число также не делится на три.

\subsection*{Почему этот шаг неверен}
В условии требуется, чтобы число делилось на три \emph{и одновременно} на четыре. Нельзя собрать варианты, удовлетворяющие только одному признаку.

Признак делимости на четыре проверяет только последние две цифры. Для числа $12\square4$ последние две цифры имеют вид $\square4$. Они делятся на четыре при цифрах
\[
0,\ 2,\ 4,\ 6,\ 8.
\]

Признак делимости на три проверяет сумму всех цифр. Здесь сумма равна
\[
1+2+\square+4=7+\square.
\]
Она делится на три при цифрах
\[
2,\ 5,\ 8.
\]
Нужны цифры, которые находятся сразу в обоих наборах. Такими цифрами являются только два и восемь.

\subsection*{Ключевая идея}
\KeyIdea\ при слове «одновременно» нужно найти пересечение условий: каждая выбранная цифра обязана пройти и признак делимости на три, и признак делимости на четыре.

\subsection*{Исправление от последнего правильного шага}
Вариант $1224$ сохраняем: он проходит обе проверки. Варианты $1204$ и $1264$ исключаем, потому что суммы их цифр не делятся на три. Вариант $1284$ сохраняем: последние две цифры $84$ делятся на четыре, а сумма цифр равна пятнадцати и делится на три.

Следовательно, остаются цифры
\[
2\quad\text{и}\quad8.
\]

\subsection*{Полное правильное решение}
Сначала применим признак делимости на четыре. Последние две цифры должны образовать число, кратное четырём:
\[
04,\ 24,\ 44,\ 64,\ 84.
\]
Значит по этому признаку возможны цифры
\[
0,\ 2,\ 4,\ 6,\ 8.
\]

Теперь применим признак делимости на три. Сумма цифр числа $12\square4$ равна $7+\square$. Проверяем цифры от нуля до девяти и получаем значения, при которых эта сумма кратна трём:
\[
\square=2,\ 5,\ 8.
\]

Пересекаем два набора:
\[
\{0,2,4,6,8\}\cap\{2,5,8\}=\{2,8\}.
\]
Значит в квадрат можно поставить только цифры два и восемь.

\subsection*{Проверка}
Для цифры два получаем $1224$: сумма цифр равна девяти, а $24$ делится на четыре.

Для цифры восемь получаем $1284$: сумма цифр равна пятнадцати, а $84$ делится на четыре.

Оба числа делятся и на три, и на четыре.

\subsection*{Итог}
\Outcome\ правильный набор цифр: $2$ и $8$. Варианты с цифрами ноль и шесть проходят проверку на делимость на четыре, но не проходят проверку на делимость на три.

\subsection*{Задача на закрепление}
В числе $13\square2$ вместо квадрата поставьте все цифры, при которых число одновременно делится на $3$ и на $4$.

\newpage
\phantomsection\label{task:3}
\section*{Задание 3}

\subsection*{Текст задания}
Разложите числа $72$ и $90$ на простые множители.

\subsection*{Ваше решение}
Для первого числа в работе начато разложение числа $70$:
\[
70:5=14,\qquad 14:2=7,\qquad 7:7=1.
\]
Для второго числа записано корректное последовательное деление:
\[
90:3=30,\qquad 30:3=10,\qquad 10:2=5,\qquad 5:5=1.
\]

\subsection*{Ваш ответ}
Для числа $72$ разложение не получено, поскольку вместо него использовано число $70$. Разложение числа $90$ по выполненным делениям соответствует множителям $2\cdot3^2\cdot5$.

\subsection*{Ошибка}
\ErrorLabel\ вместо числа $72$ разложено число $70$.

\subsection*{Где именно возникает ошибка}
Для ветви с числом $72$ последнего вычислительного шага до ошибки нет: ошибка возникает уже в первом переписывании исходного числа. В условии дано $72$, а в столбике записано $70$. После этого деления на пять, два и семь могут быть арифметически согласованы с числом семьдесят, но они уже не решают поставленную задачу.

Ветка с числом $90$ выполнена по существу верно, поэтому её можно сохранить.

\subsection*{Почему этот шаг неверен}
Разложение на простые множители --- это равенство между исходным числом и произведением простых чисел. Если исходное число изменить, получится разложение другого числа.

Записанные множители для первой ветви дают
\[
5\cdot2\cdot7=70,
\]
то есть подтверждают именно число семьдесят. Но нужно получить произведение, равное семидесяти двум. Простые множители числа семьдесят два другие: в нём нет множителей пять и семь.

\subsection*{Ключевая идея}
\KeyIdea\ в лестнице разложения нужно начинать строго с данного числа и на каждом шаге делить его на простой делитель без остатка, пока не получится единица.

\subsection*{Исправление от последнего правильного шага}
Для числа $90$ уже записанный ход можно оставить. Исправить нужно первую ветвь: вместо $70$ начинаем с $72$.
\[
72:2=36,\qquad
36:2=18,\qquad
18:2=9,\qquad
9:3=3,\qquad
3:3=1.
\]
Отсюда
\[
72=2\cdot2\cdot2\cdot3\cdot3=2^3\cdot3^2.
\]

\subsection*{Полное правильное решение}
Разложим $72$. Делим последовательно на простые числа:
\[
72:2=36,
\]
\[
36:2=18,
\]
\[
18:2=9,
\]
\[
9:3=3,
\]
\[
3:3=1.
\]
Следовательно,
\[
72=2\cdot2\cdot2\cdot3\cdot3=2^3\cdot3^2.
\]

Теперь разложим $90$. Можно сохранить порядок делений из работы:
\[
90:3=30,
\]
\[
30:3=10,
\]
\[
10:2=5,
\]
\[
5:5=1.
\]
Следовательно,
\[
90=3\cdot3\cdot2\cdot5=2\cdot3^2\cdot5.
\]

\subsection*{Проверка}
Проверим произведения:
\[
2^3\cdot3^2=8\cdot9=72,
\]
\[
2\cdot3^2\cdot5=2\cdot9\cdot5=90.
\]
Оба произведения возвращают исходные числа.

\subsection*{Итог}
\Outcome\ для числа $90$ ход решения можно сохранить. В первой ветви нужно заменить ошибочно переписанное число $70$ на исходное $72$ и получить разложение $2^3\cdot3^2$.

\subsection*{Задача на закрепление}
Разложите числа $84$ и $126$ на простые множители.

\newpage
\phantomsection\label{task:5}
\section*{Задание 5}

\subsection*{Текст задания}
Найдите $\text{НОД}(414,504)$, затем сократите дробь
\[
\frac{414}{504}.
\]

\subsection*{Ваше решение}
Разложения на простые множители записаны верно:
\[
414=2\cdot3\cdot3\cdot23=2\cdot3^2\cdot23,
\]
\[
504=2\cdot2\cdot2\cdot3\cdot3\cdot7=2^3\cdot3^2\cdot7.
\]
После этого в работе записано
\[
\text{НОД}(414,504)=2\cdot3\cdot7=42.
\]
Сокращённая дробь не записана.

\subsection*{Ваш ответ}
$\text{НОД}(414,504)=42$; требуемое сокращение дроби не завершено.

\subsection*{Ошибка}
\ErrorLabel\ неверно выбраны общие простые множители для НОД; из-за этого не выполнена вторая часть задания.

\subsection*{Где именно возникает ошибка}
Последний полностью правильный шаг --- два разложения
\[
414=2\cdot3^2\cdot23,
\qquad
504=2^3\cdot3^2\cdot7.
\]
Первый неверный шаг --- переход к произведению
\[
2\cdot3\cdot7.
\]
Здесь множитель семь взят ошибочно: семь есть в разложении числа $504$, но отсутствует в разложении числа $414$. Кроме того, число три является общим множителем два раза, потому что в обоих разложениях присутствует $3^2$, однако в записанный НОД вошла только одна тройка.

\subsection*{Почему этот шаг неверен}
НОД по разложению на простые множители составляют только из тех простых множителей, которые есть \emph{в обоих} числах. Для каждого общего простого множителя берут меньшую из двух степеней.

Для двойки степени равны один и три, поэтому в НОД входит одна двойка. Для тройки степени равны два и два, поэтому в НОД входят две тройки. Множитель семь есть только во втором числе, поэтому он в НОД не входит. Множитель двадцать три есть только в первом числе и также не входит.

Это даёт
\[
\text{НОД}(414,504)=2\cdot3^2=18.
\]
Число $42$ не может быть НОД ещё и по простой проверке: оно должно делить оба исходных числа без остатка, а $414$ на $42$ без остатка не делится.

\subsection*{Ключевая идея}
\KeyIdea\ после правильного разложения нужно не собирать все встреченные простые числа, а взять пересечение разложений с наименьшими степенями. Затем найденный НОД удобно использовать для сокращения дроби сразу за один шаг.

\subsection*{Исправление от последнего правильного шага}
Сохраняем правильные разложения и выбираем общие множители:
\[
414=2^1\cdot3^2\cdot23,
\qquad
504=2^3\cdot3^2\cdot7.
\]
Общая часть равна
\[
2^1\cdot3^2=18.
\]
Значит
\[
\text{НОД}(414,504)=18.
\]
Теперь сокращаем дробь на восемнадцать:
\[
\frac{414}{504}
=\frac{414:18}{504:18}
=\frac{23}{28}.
\]

\subsection*{Полное правильное решение}
Разложим числа на простые множители:
\[
414=2\cdot207=2\cdot3\cdot69=2\cdot3\cdot3\cdot23
=2\cdot3^2\cdot23,
\]
\[
504=2\cdot252
=2^2\cdot126
=2^3\cdot63
=2^3\cdot3\cdot21
=2^3\cdot3^2\cdot7.
\]

Сравниваем степени общих простых множителей. Для двойки берём степень один, для тройки --- степень два:
\[
\text{НОД}(414,504)=2\cdot3^2=18.
\]

Делим числитель и знаменатель дроби на их НОД:
\[
414:18=23,
\qquad
504:18=28.
\]
Поэтому
\[
\frac{414}{504}=\frac{23}{28}.
\]

\subsection*{Проверка}
Проверяем НОД:
\[
18\cdot23=414,
\qquad
18\cdot28=504.
\]
Значит восемнадцать действительно делит оба числа.

После сокращения числа $23$ и $28$ не имеют общего делителя больше единицы: двадцать три --- простое число и не делит двадцать восемь. Следовательно, дробь $\dfrac{23}{28}$ уже несократима.

\subsection*{Итог}
\Outcome\ правильный НОД равен $18$, а исходная дробь сокращается до $\dfrac{23}{28}$. Ошибка появилась именно при выборе общих простых множителей, хотя сами разложения двух чисел были записаны верно.

\subsection*{Задача на закрепление}
Найдите $\text{НОД}(270,378)$ по разложениям на простые множители и сократите дробь
\[
\frac{270}{378}.
\]

\end{document}
